\chapter{Experimental Design}
\label{ch:experiments}

Four experiments are conducted to evaluate the tau-augmented agent against the baseline.

\paragraph{Experiment A -- Learning curves.}
Both agents are trained for $1\,000\,000$ timesteps across three random seeds (0, 1, 2).
The mean touchdown speed $|v_z|$ and success rate are logged every $20\,000$ timesteps
to compare convergence speed and training stability.

\paragraph{Experiment B -- Performance comparison.}
After training, both agents are evaluated on 30 \emph{paired} episodes per seed
(90 episodes total per agent).
Paired evaluation means both agents are initialised with identical start conditions
(controlled by a fixed episode seed), enabling a direct, within-episode comparison.
Two metrics are collected per episode:
\begin{itemize}
  \item Touchdown speed $|v_z|$ at the moment of landing.
  \item Tau-regulation error: the mean $|\dot{\tau} - (-0.5)|$ computed across all
        descent steps of that episode.
\end{itemize}
Statistical significance is assessed using the Wilcoxon signed-rank test (paired,
two-tailed, $n = 90$) and bootstrap 95\% confidence intervals
($n_\text{boot} = 2000$).

\paragraph{Experiment C -- Wind robustness (zero-shot).}
Both agents are evaluated under horizontal wind disturbances at five levels
$w \in \{0.0,\, 0.5,\, 1.0,\, 1.5,\, 2.0\}$~m/s$^2$ (see Section~\ref{sec:wind_model}).
Fifty episodes per seed per wind level are run and success rate and mean $|v_z|$ are
reported across seeds.
Neither agent was trained with wind, so this tests zero-shot generalisation.

\paragraph{Experiment D -- Sensitivity analysis.}
The effect of the shaping coefficient $\lambda$ is investigated by training the tau
agent at six values:
$\lambda \in \{0.00,\, 0.02,\, 0.05,\, 0.10,\, 0.20,\, 0.50\}$.
All runs use seed 0 and $500\,000$ training steps (half the main budget) to keep
computation manageable; this is noted as a limitation in
Chapter~\ref{ch:conclusion}.
Performance is evaluated on 30 episodes with mean $|v_z|$ and tau-regulation error
as metrics.
